% Sample Applications, Chapter 7 -- Price List (annex Sec 8.3).
% Values replayed through the shipped application (examples/prices7.i23);
% @ROUND as os/i123/core/eval.c computes it, Percent and Currency as
% os/i123/core/fmt.c.  Figures screens/sa7.k123, sa7b.k123.
% GATES: test-i123-formula test-pubs-i123-examples test-pubs-screens
\Gates{test-i123-formula, test-pubs-i123-examples, test-pubs-screens}
\chapter{Price List}

When costs rise, the shop raises every price by the same percentage and
rounds each new price to the cent.  This worksheet keeps the percentage in
one cell, so a single entry reprices the whole list.

\section{What You Enter}
\begin{itemize}
\item The percent increase
\item The name and present price of each item
\end{itemize}

\section{Entering the Data}
\ix{price list}\ix{"@ROUND}
Type \type{Rate} in A1 and the increase, \type{.05} for 5 percent, in B1;
format B1 with \slashcmd{/Range Format Percent}, no decimal places.  Type the
headings \type{Item}, \type{"Old} and \type{"New} in A3..C3, the items
\type{Pen}, \type{Pad} and \type{Tape} in A4..A6 and the present prices
\type{1.5}, \type{2.25} and \type{.8} in B4..B6.  In C4 enter
\type{@round(b4*(1+\$b\$1),2)} and copy it to C5 and C6.  Format B4..C6 as
currency with two decimal places.
\FlairShot[fig:sa7]{sa7}{Present and new prices}

\begin{LotTable}{p{30mm}p{24mm}p{70mm}}
\LotHead{\textbf{Item} & \textbf{Address} & \textbf{Meaning}}
Rate & B1 & The percent increase\\
Item & A4..A6 & The name of each item\\
Old & B4..B6 & The present price\\
New & C4..C6 & The price after the increase, rounded to the cent\LotEnd
\end{LotTable}

\section{Interpreting the Results}
\LotExample{examples/prices7}
The formula multiplies the old price by one plus the rate, and
\fn{@ROUND} keeps two decimal places.  Without it a cell would hold a price
such as a fraction of a cent, which displays rounded but still carries the
extra digits into later formulas.

\section{Changing the Template}
To try a different increase, type a new rate in B1; the three new prices
follow at once.  Figure~\ref{fig:sa7b} shows the prices at an increase of
8 percent, typed as \type{.08}.
\FlairShot[fig:sa7b]{sa7b}{New prices at an increase of 8 percent}
  To add an item, type it below Tape and copy the formula in
C6 to the new row.
